\section*{الفصل \ref{ch:b2:living-soil} --- التربة الحية}

\begin{solution}{exo:b2:living-soil:1}
O: فُتات ومادة عضوية متحللة جزئيًا، وفيه يبدأ
التحلل. وA: \omterm{def:b2:living-soil:soil}{تربة} معدنية ممتزجة بالدُّبال، داكنة ذات بنية فُتاتية،
تحمل معظم الجذور والكائنات والمغذيات والماء. وB: أفق
التراكم، حيث يترسب الطين وأكاسيد الحديد والكربونات المغسولة
من فوق. وC: مادة أم مجوّاة، وشظايا صخر
في حشوة ناعمة، وفيها تُطلق \omterm{def:b2:living-soil:soil}{تربة} معدنية جديدة. وR:
صخر قاع غير مجوّى.
\end{solution}

\begin{solution}{exo:b2:living-soil:2}
الشحنات السالبة على سطوح الطين والدُّبال تمسك شوارد موجبة
قابلة للتبادل يستطيع الجذر التقاطها مقابل بروتونات. أما حبيبات الرمل
فمرو، غير مشحونة وخشنة، وسطحها في الغرام قليل
ولا \omterm{def:b2:living-soil:soil}{دُبال} فيها يُذكر. والتكليس يضيف كربونات الكالسيوم:
فيزيح الكالسيوم شوارد الهدروجين والألمنيوم التي تشبع
مواقع التبادل في \omterm{def:b2:living-soil:soil}{تربة} حمضية، وتعادل الكربونات
الحموضة، وتُملأ المواقع من جديد بشاردة مغذية.
\end{solution}

\begin{solution}{exo:b2:living-soil:3}
المحلِّلات: البكتيريا (\emph{Bacillus}) والفطريات
(\emph{Trichoderma})؛ وراعيات المجهريات: الأوّليّات (الأميبات)
والخيطيات آكلة البكتيريا؛ وآكلات الحطام: ديدان الأرض وأم أربعة وأربعين
والقافزات الذنبية وقمل \omterm{def:b2:plant-meristems:wood}{الخشب}؛ والمفترسات: العث والخيطيات المفترسة
وأم أربعة وأربعين المفترسة وخنافس الأرض؛ والخُلد أو الزبابة في القمة. و
الراعيات، و$\mathrm{C:N}$ عندها أعلى منها عند فرائسها البكتيرية،
تطرح فائض النتروجين أمونيومًا: فهي تعدّن
النتروجين الذي حبسته البكتيريا.
\end{solution}

\begin{solution}{exo:b2:living-soil:4}
الشجيرية: يدخل الفطر (وهو غلوميرومِسِيت) خلايا قشرة
الجذر ويكوّن شُجيرات؛ وتوجد في معظم الأعشاب والمحاصيل وفي
أشجار المناطق المدارية. والخارجية: يغمد الفطر (وهو من البازيديات
والزقيات، وكثير منها فطر مظلي) طرف الجذر
وينمو بين الخلايا دون أن يدخلها؛ وتوجد في أشجار المناطق المعتدلة
والشمالية. وفي الحالتين يعطي النبات سكرًا ويعطي الفطر
فوسفاتًا ونتروجينًا وماءً ومغذيات دقيقة وشيئًا من
الحماية من الممرضات.
\end{solution}

\begin{solution}{exo:b2:living-soil:5}
$L^{*} = I/k = 4/0.8 = \qty{5}{t/ha}$؛ \omterm{def:b2:biogeochemical-cycles:reservoir}{وزمن المكث} $1/k = 1.25$
سنة؛ وفقد \qty{90}{\%} بعد $\ln 10/0.8 = 2.9$ سنة.
\end{solution}

\begin{solution}{exo:b2:living-soil:6}
لكل \qty{100}{kg} من الكربون المأكول تحتفظ المجهريات بمقدار \qty{40}{kg} و
تحتاج $40/8 = \qty{5}{kg}$ من النتروجين. والقش يمد $100/80 =
\qty{1.25}{kg}$: أي تثبيط قدره \qty{3.75}{kg} لكل \qty{100}{kg}
من الكربون. والبرسيم يمد $100/15 = \qty{6.7}{kg}$: أي تمعدن
قدره \qty{1.7}{kg}.
\end{solution}

\begin{solution}{exo:b2:living-soil:7}
\qty{20}{t/ha} من \omterm{def:b2:living-soil:soil}{التربة} في السنة عبر الديدان. و\qty{20}{cm}
العليا تزن $0.2\times 10^{4}\times 1.3 = \qty{2600}{t/ha}$:
أي 130 سنة لمرور واحد. وعشرة أطنان داروين في الفدان هي
\qty{25}{t/ha}، أي الرتبة نفسها --- و«كل بضع سنين» عنده تشير
إلى الدُّبال السطحي، أي بضعة سنتيمترات، لا إلى \omterm{def:b2:living-soil:soil}{التربة} السطحية كلها.
\end{solution}

\begin{solution}{exo:b2:living-soil:8}
البكتيريا: $10^{9}\times 10^{-12} = \qty{1}{mg}$ في الغرام. والهيفات:
حجمها $\pi(2\times 10^{-4})^{2}\times 10^{4} = \qty{1.26e-3}{cm^3}$،
وكتلتها \qty{1.4}{mg} في الغرام. وفي الهكتار ($2.6\times 10^{9}$ غرام):
\qty{2.6}{t} من البكتيريا و\qty{3.6}{t} من الفطريات.
\end{solution}

\begin{solution}{exo:b2:living-soil:9}
$H^{*} = 1.5/0.02 = \qty{75}{t/ha}$ من الكربون. وبالحراثة: $1.5/0.04
= \qty{37.5}{t/ha}$؛ والفقد \qty{37.5}{t/ha}، ويزول نصفه بعد
$\ln 2/0.04 = 17$ سنة.
\end{solution}

\begin{solution}{exo:b2:living-soil:10}
الفوسفات: $\sqrt{2\times 10^{-13}\times 8.64\times 10^{6}} =
\qty{1.3}{mm}$ في الموسم --- فالجذر الذي يبعد \qty{1}{cm} عن التالي
يبلغ ثُمن \omterm{def:b2:living-soil:soil}{التربة} بينهما، في حين تبلغ الهيفات المتباعدة
\qty{0.1}{mm} كلها. والنترات: \qty{4}{cm}، أي أبعد
من تباعد الجذور: فهي تأتي إلى الجذر بنفسها (وبالجريان
الكتلي مع الماء)، فتضيف الميكوريزا قليلًا للنترات
وكثيرًا للفوسفات.
\end{solution}

\begin{solution}{exo:b2:living-soil:11}
الكتلة $0.25\times 10^{4}\times 1.3 = \qty{3250}{t/ha}$. والفقد الصافي
\qty{19}{t/ha} في السنة: فالعمر 170 سنة. والكربون المفقود $20\times
0.03 = \qty{0.6}{t/ha}$ في السنة. وتحت الزراعة دون حراثة يكون الفقد الصافي
\qty{1}{t/ha}: أي 3250 سنة.
\end{solution}

\begin{solution}{exo:b2:living-soil:12}
الفُتات يدور في سنة أو سنتين، والدُّبال في خمسين، والأفق A
المعدني في آلاف (إذ يتكون بمعدل \qty{0.1}{mm} في السنة)، و
الملح يتراكم في عقود ويُصرف في عقود إن كان هناك
صرف. والمزارع يرى الفُتات والمحصول يستجيبان في غضون
موسم ويرى الدُّبال في غضون مسيرة عمره (فهبوط بمقدار الثلث يستغرق
ثلاثين سنة؛ والتعافي مثلها)؛ أما \omterm{prop:b2:living-soil:erosion}{تعرية التربة} المعدنية
فخفية في عمر واحد ولا تنعكس في عشرة؛ والملح
مرئي في جيل، ودائم دون صرف. أما
المتغيرات البطيئة في \omterm{def:b2:living-soil:soil}{التربة} فلا يسجلها أي حساب سنوي.
\end{solution}

\begin{solution}{pb:b2:living-soil:1}
\textbf{1.} $0.25\times 10^{4}\times 1.3 = \qty{3250}{t/ha}$؛
والكربون $\qty{81}{t/ha}$.
\textbf{2.} $81/12 = \qty{6.8}{t/ha}$ من النتروجين.
\textbf{3.} البكتيريا \qty{1}{mg/g}، أي $\qty{3.25}{t/ha}$؛ والهيفات
$\pi(2\times 10^{-4})^{2}\times 2\times 10^{4}\times 1.1 =
\qty{2.8}{mg/g}$، أي \qty{9}{t/ha}: فالكتلة الحيوية المجهرية نحو
\qty{12}{t/ha}.
\textbf{4.} نحو \qty{14}{t/ha} من المادة الحية (بالكتلة الطازجة)،
أي نحو \qty{17}{\%} من كربون الدُّبال --- وبضعة في المئة منه
كربونًا، إذ الكائنات ماء في معظمها.
\textbf{5.} الكربون المجهري نحو نصف \qty{12}{t}، أي \qty{6}{t}؛
وإذ تدور مرتين في السنة مع الاحتفاظ بمقدار \qty{40}{\%}، تستهلك
$2\times 6/0.4 = \qty{30}{t}$ من الكربون وتتنفس \qty{18}{t/ha}
من الكربون، أي \qty{66}{t} من $\mathrm{CO_2}$ (وهو تقدير خشن: فالاحتفاظ
والدوران تقريبيان).
\textbf{6.} ثلاثة أمثال الإنتاج الأولي الصافي للمحصول البالغ
\qty{5}{t}: فالتقدير مرتفع أكثر مما ينبغي، إذ الفرضيات (دوران الكتلة
الحيوية كلها مرتين، ولا شيء منها ساكن) سخية؛
\omterm{prop:b2:living-soil:humus}{وتنفس التربة} الحقيقي تحت محصول من رتبة إنتاج المحصول
نفسه، أي عدة أطنان من الكربون في الهكتار في السنة. \omterm{def:b2:living-soil:soil}{فالتربة}
تتنفس بقدر ما تتنفس النباتات فوقها.
\textbf{7.} الكربون $5\times 0.45 = \qty{2.25}{t}$؛ والنتروجين $5\times
0.03 = \qty{150}{kg}$؛ و$\mathrm{C:N} = 15$.
\textbf{8.} تُمعدِن: فالمجهريات تحتاج $2250\times 0.4/8 =
\qty{112}{kg}$ من النتروجين والبقايا تحمل 150، فيُطلق
\qty{38}{kg/ha} صافيًا.
\textbf{9.} $1 - e^{-0.5} = \qty{39}{\%}$ بعد 3 أشهر؛ و$1 - e^{-2}
= \qty{86}{\%}$ بعد سنة.
\textbf{10.} $0.39\times 38 = \qty{15}{kg/ha}$.
\textbf{11.} يحتاج المحصول \qty{150}{kg}، وجُلّها في شهريه
الأولين؛ والبقايا تعطي \qty{15}{kg} في ثلاثة أشهر و
\qty{38}{kg} إجمالًا --- أي مكمّلًا لا موردًا. وقيمتها
أيضًا في النتروجين الذي ثبّته البرسيم ويحمله في
\qty{150}{kg} البقايا نفسها، ويدخل جلّه الدُّبال ويعود
على مدى سنين.
\textbf{12.} القش: الكربون \qty{2250}{kg}، والنتروجين $2250/80 =
\qty{28}{kg}$؛ والمجهريات تحتاج \qty{112}{kg}، فيُثبَّط \qty{84}{kg/ha}
من \omterm{def:b2:living-soil:soil}{التربة} --- أي أكثر من نصف حاجة المحصول.
وعلى المزارع أن يضيف نتروجينًا مع القش، أو يحرثه قبل
البذر بوقت كافٍ، أو يتركه على السطح حيث يتحلل ببطء.
\textbf{13.} بالحراثة: $1.2/0.025 = \qty{48}{t/ha}$؛ ودون حراثة:
$1.2/0.015 = \qty{80}{t/ha}$.
\textbf{14.} $H(t) = 80 - 32\,e^{-0.015t}$: أي كسب \qty{4.5}{t/ha}
بعد 10 سنوات، و\qty{17}{t/ha} بعد 50.
\textbf{15.} $\mathrm{d}H/\mathrm{d}t = 0.015\times 32 =
\qty{0.48}{t/ha}$ من الكربون في السنة الأولى، أي \qty{1.8}{t} من
$\mathrm{CO_2}$.
\textbf{16.} $0.48\times 1.5\times 10^{9} = \qty{0.7}{GtC}$ في السنة
الأولى، أي \qty{7}{\%} من الانطلاقات الأحفورية؛ وبحلول السنة الخمسين
يكون المعدل قد هبط إلى $0.48\,e^{-0.75} = \qty{0.23}{t/ha}$،
أي \qty{0.34}{GtC}.
\textbf{17.} الكسب هو الاقتراب من \omterm{def:b2:biogeochemical-cycles:reservoir}{حالة مستقرة} جديدة: فهو
يخبو إلى الصفر إذ يبلغ $H$ \qty{80}{t/ha}، والكسب كله
يُعاد إلى الهواء في غضون عقود إذا حُرث الحقل من جديد
--- فمصرف \omterm{def:b2:living-soil:soil}{التربة} نقل مرة واحدة قابل للانعكاس، لا تعويض دائم
عن انطلاقات مستمرة.
\textbf{18.} عند $k_h = 0.02$: يكون $H^{*} = \qty{60}{t/ha}$، أي فقد
\qty{20}{t/ha}، وهو أكثر من كسب خمسين سنة من الزراعة دون حراثة:
فالاحترار يستطيع نقض
الإدارة.
\textbf{19.} بالحراثة: $15 - 0.5 = \qty{14.5}{t/ha}$ في السنة؛ ودون حراثة:
\qty{1}{t/ha}.
\textbf{20.} $3250/14.5 = 220$ سنة؛ و$3250/1 = 3250$ سنة.
\textbf{21.} الكربون $15\times 0.025 = \qty{0.38}{t/ha}$؛ والنتروجين
$0.38/12 = \qty{31}{kg/ha}$ في السنة.
\textbf{22.} خبو الدُّبال عند \omterm{def:b2:biogeochemical-cycles:reservoir}{الحالة المستقرة} المحروثة هو $k_hH^{*} =
I = \qty{1.2}{t/ha}$ في السنة؛ وتضيف التعرية ثلثًا آخر،
لكنها، بخلاف الخبو، غير موازَنة بدخل --- فهي تستنزف
المخزون.
\textbf{23.} جزئيًا. فالكربون المتعري المدفون في رسوب أو
خزان قد يكون محميًا من الخبو زمنًا طويلًا (أي مصرف
دفن)؛ لكن التجمعات تُكسر في النقل، ويتأكسد جل
الكربون في الطريق، وخصوبة الحقل المفقودة
يُستعاض عنها بسماد ينطلق في صنعه كربون. وإشارة
أثر التعرية في الكربون عالميًا ما تزال موضع جدل.
\textbf{24.} الفُتات $1/k = 0.5$ سنة؛ والدُّبال $1/k_h = 40$ سنة
بالحراثة و67 دون حراثة؛ \omterm{def:b2:living-soil:soil}{والتربة} المعدنية $3250/0.5 =
6500$ سنة بالنسبة إلى التكوّن، و220 سنة بالنسبة
إلى الفقد بالحراثة.
\textbf{25.} مخزون الكربون \qty{81}{t/ha} في الأفق A؛ والعمر
220 سنة بالحراثة، و3250 دون حراثة؛ وبقايا البرسيم
تمد نحو \qty{38}{kg/ha} من النتروجين المعدني، منها 15 في
الأشهر الثلاثة الأولى.
\end{solution}
